\documentclass[11pt]{article}
\usepackage[a4paper,margin=1in]{geometry}
\usepackage{fontspec}
\usepackage[boldfont=false]{xeCJK}
\setCJKmainfont{FandolSong-Regular.otf}
\usepackage{amsmath}
\usepackage{amssymb}
\usepackage{array}
\usepackage{booktabs}
\usepackage{graphicx}
\usepackage{adjustbox}
\usepackage{enumitem}
\setlist[itemize]{topsep=2pt,partopsep=0pt,parsep=0pt,itemsep=2pt,leftmargin=1.4em}
\setlist[enumerate]{topsep=2pt,partopsep=0pt,parsep=0pt,itemsep=2pt,leftmargin=1.6em}
\usepackage{parskip}
\usepackage{fvextra}
\fvset{breaklines=true,breakanywhere=true,fontsize=\small}
\setlength{\parindent}{0pt}

\begin{document}

\begin{center}
  {\LARGE\bfseries Algebra and graphs}\\[0.35em]
  {\large Igcse Mathematics}
\end{center}
\vspace{0.6em}

This handout covers Topic 2, Algebra and graphs. Parts marked \textbf{(Extended)} are only tested on the Extended papers; everything else is for both levels.

\section*{Working with algebra}

In \textbf{algebra} 代数 we use letters to stand for numbers. A letter whose value can change is a \textbf{variable} 变量. To \textbf{substitute} 代入 means to put a number in place of a letter.

\textbf{Worked example.} Find the value of $3x^{2} - 2y$ when $x = 4$ and $y = 5$.

\[
3 \times 4^{2} - 2 \times 5 = 3 \times 16 - 10 = 48 - 10 = 38.
\]

\section*{Simplifying and expanding}

A \textbf{term} 项 is a single part of an \textbf{expression} 表达式, such as $5a$ or $-9b$. \textbf{Like terms} 同类项 have exactly the same letters; you may add or subtract them. The number in front of the letter is the \textbf{coefficient} 系数.

\textbf{Worked example.} Simplify $2a^{2} + 3ab - 1 + 5a^{2} - 9ab + 4$.

Collect like terms: $2a^{2} + 5a^{2} = 7a^{2}$, $3ab - 9ab = -6ab$, $-1 + 4 = 3$. So the answer is

\[
7a^{2} - 6ab + 3.
\]

To \textbf{expand} 展开 means to multiply out \textbf{brackets} 括号. Multiply every term inside by the term outside; for two brackets, multiply every term in the first by every term in the second.

\textbf{Worked examples.}

\[
3x(2x - 4y) = 6x^{2} - 12xy.
\]

\[
(2x + 1)(x - 4) = 2x^{2} - 8x + x - 4 = 2x^{2} - 7x - 4.
\]

For three brackets (\textbf{Extended}), expand two first, then multiply by the third:

\[
(x - 2)(x + 3)(2x + 1) = (x^{2} + x - 6)(2x + 1) = 2x^{3} + 3x^{2} - 11x - 6.
\]

\section*{Factorising}

To \textbf{factorise} 因式分解 is the opposite of expanding: write the expression as a product of brackets. Always take out the \textbf{common factor} 公因式 first.

\textbf{Worked example.} $9x^{2} + 15xy = 3x(3x + 5y)$, because $3x$ divides both terms.

The following patterns are \textbf{Extended}.

\textbf{Grouping} (four terms): take a common factor from each pair.

\[
xy + 2x + 3y + 6 = x(y + 2) + 3(y + 2) = (x + 3)(y + 2).
\]

\textbf{Difference of two squares} 平方差: $a^{2} - b^{2} = (a + b)(a - b)$.

\[
9x^{2} - 16 = (3x + 4)(3x - 4).
\]

\textbf{Perfect square} 完全平方: $a^{2} + 2ab + b^{2} = (a + b)^{2}$.

\[
x^{2} + 6x + 9 = (x + 3)^{2}.
\]

\textbf{Quadratic} 二次 expressions $ax^{2} + bx + c$: find two numbers that multiply to $a \times c$ and add to $b$, then split the middle term.

\textbf{Worked example.} Factorise $2x^{2} + 7x + 3$.

Here $a \times c = 6$ and $b = 7$. The numbers $1$ and $6$ work. Split and group:

\[
2x^{2} + x + 6x + 3 = x(2x + 1) + 3(2x + 1) = (x + 3)(2x + 1).
\]

For $ax^{3} + bx^{2} + cx$, take out the common $x$ first: $2x^{3} + 7x^{2} + 3x = x(2x^{2} + 7x + 3) = x(x + 3)(2x + 1)$.

\section*{Completing the square (Extended)}

\textbf{Completing the square} 配方法 rewrites $x^{2} + bx + c$ as $(x + p)^{2} + q$. Take half of the $x$-coefficient, square it, then balance.

\textbf{Worked example.} Write $x^{2} + 6x + 1$ in completed-square form.

Half of $6$ is $3$, and $3^{2} = 9$:

\[
x^{2} + 6x + 1 = (x + 3)^{2} - 9 + 1 = (x + 3)^{2} - 8.
\]

When the $x^{2}$ has a coefficient, take it out of the first two terms first:

\[
2x^{2} + 8x + 3 = 2(x^{2} + 4x) + 3 = 2\big((x + 2)^{2} - 4\big) + 3 = 2(x + 2)^{2} - 5.
\]

\section*{Algebraic fractions (Extended)}

An \textbf{algebraic fraction} 分式 has algebra on the top or bottom. Add and subtract using a common denominator; multiply and divide as with ordinary fractions.

\textbf{Worked examples.}

\[
\frac{x}{3} + \frac{x - 4}{2} = \frac{2x}{6} + \frac{3(x - 4)}{6} = \frac{2x + 3x - 12}{6} = \frac{5x - 12}{6}.
\]

\[
\frac{3a}{4} \div \frac{9a}{10} = \frac{3a}{4} \times \frac{10}{9a} = \frac{30a}{36a} = \frac{5}{6}.
\]

To simplify a \textbf{rational expression} 有理式, factorise the top and bottom, then cancel common brackets.

\[
\frac{x^{2} - 2x}{x^{2} - 5x + 6} = \frac{x(x - 2)}{(x - 2)(x - 3)} = \frac{x}{x - 3}.
\]

\section*{Indices in algebra}

The laws of \textbf{indices} 指数 work with letters too: $a^{m} \times a^{n} = a^{m+n}$, $a^{m} \div a^{n} = a^{m-n}$, and $(a^{m})^{n} = a^{mn}$.

\textbf{Worked examples.}

\[
(5x^{3})^{2} = 25x^{6}, \qquad 12a^{5} \div 3a^{-2} = 4a^{7}, \qquad 6x^{7}y^{4} \times 5x^{-5}y = 30x^{2}y^{5}.
\]

You can also solve simple index equations by writing both sides with the same \textbf{base} 底数.

\textbf{Worked example.} Solve $2^{x} = 32$. Since $32 = 2^{5}$, you get $x = 5$.

\section*{Equations}

An \textbf{equation} 方程 says two expressions are equal. To solve a \textbf{linear} 一次 equation, do the same operation to both sides until the \textbf{unknown} 未知数 is alone.

\textbf{Worked example.} Solve $5 - 2x = 3(x + 7)$.

\[
5 - 2x = 3x + 21 \;\Rightarrow\; 5 - 21 = 3x + 2x \;\Rightarrow\; -16 = 5x \;\Rightarrow\; x = -\tfrac{16}{5}.
\]

\subsection*{Fractional equations (Extended)}

A \textbf{fractional equation} 分式方程 has the unknown in a denominator. Multiply both sides by the denominator to clear it.

\textbf{Worked example.} Solve $\dfrac{x}{2x + 1} = 4$.

\[
x = 4(2x + 1) = 8x + 4 \;\Rightarrow\; -7x = 4 \;\Rightarrow\; x = -\tfrac{4}{7}.
\]

\subsection*{Simultaneous equations}

\textbf{Simultaneous equations} 联立方程 are two equations solved together. For two linear equations, add or subtract to remove one letter.

\textbf{Worked example.} Solve $2x + y = 7$ and $3x - y = 8$.

Adding removes $y$: $5x = 15$, so $x = 3$. Then $y = 7 - 2(3) = 1$.

For one linear and one \textbf{quadratic} equation (\textbf{Extended}), substitute the linear into the curve.

\textbf{Worked example.} Solve $y = x + 2$ and $y = x^{2}$.

\[
x^{2} = x + 2 \;\Rightarrow\; x^{2} - x - 2 = 0 \;\Rightarrow\; (x - 2)(x + 1) = 0,
\]

so $x = 2$ (giving $y = 4$) or $x = -1$ (giving $y = 1$).

\subsection*{Solving quadratic equations (Extended)}

There are three methods.

\begin{itemize}
\item \textbf{By factorising:} $x^{2} + 5x + 6 = 0 \Rightarrow (x + 2)(x + 3) = 0 \Rightarrow x = -2$ or $x = -3$.

\item \textbf{By completing the square:} $x^{2} + 6x + 1 = 0 \Rightarrow (x + 3)^{2} = 8 \Rightarrow x + 3 = \pm 2\sqrt{2} \Rightarrow x = -3 \pm 2\sqrt{2}$.

\item \textbf{By the quadratic formula} 求根公式, which is given in the exam: \[x = \frac{-b \pm \sqrt{b^{2} - 4ac}}{2a}.\]

\end{itemize}

\textbf{Worked example (formula).} Solve $2x^{2} + 3x - 1 = 0$. Here $a = 2$, $b = 3$, $c = -1$:

\[
x = \frac{-3 \pm \sqrt{9 + 8}}{4} = \frac{-3 \pm \sqrt{17}}{4}.
\]

\subsection*{Changing the subject}

To \textbf{change the subject} 公式变形 of a formula means to rearrange it so a chosen letter is alone on one side.

\textbf{Worked example.} Make $r$ the subject of $A = \pi r^{2}$ (\textbf{Extended}, because of the power).

\[
r^{2} = \frac{A}{\pi} \;\Rightarrow\; r = \sqrt{\frac{A}{\pi}}.
\]

When the letter appears twice (\textbf{Extended}), collect those terms and factorise. To make $x$ the subject of $y = \dfrac{x + 1}{x - 1}$:

\[
y(x - 1) = x + 1 \;\Rightarrow\; yx - x = 1 + y \;\Rightarrow\; x(y - 1) = 1 + y \;\Rightarrow\; x = \frac{1 + y}{y - 1}.
\]

\section*{Inequalities}

An \textbf{inequality} 不等式 uses $<$, $>$, $\leqslant$ or $\geqslant$. Solve it like an equation, but \textbf{reverse the sign} if you multiply or divide by a negative number.

\textbf{Worked example.} Solve $-3 \leqslant 3x - 2 < 7$.

Add $2$ to all parts, then divide by $3$:

\[
-1 \leqslant 3x < 9 \;\Rightarrow\; -\tfrac{1}{3} \leqslant x < 3.
\]

On a \textbf{number line} 数轴, use an open circle for $<$ or $>$ (value not included) and a closed circle for $\leqslant$ or $\geqslant$ (value included).

\subsection*{Regions (Extended)}

An inequality in two letters describes a \textbf{region} 区域 of the graph. Draw the boundary line (broken for $<$ or $>$, solid for $\leqslant$ or $\geqslant$) and shade the unwanted side. You may also be asked to list the inequalities that define a given region.

\section*{Sequences}

A \textbf{sequence} 数列 is a list of numbers that follow a rule. The \textbf{term-to-term rule} 递推规则 tells you how to get the next term from the one before.

To find the rule for the term in position $n$ (the $n$th \textbf{term}), look at how the terms change.

\textbf{Linear sequence} (the terms go up by the same amount each time). The difference is the multiple of $n$.

\textbf{Worked example.} Find the $n$th term of $2, 5, 8, 11, \dots$

The terms go up by $3$, so start with $3n$. Since $3 \times 1 = 3$ but the first term is $2$, subtract $1$: the $n$th term is $3n - 1$.

\textbf{Quadratic sequence} (the differences themselves change by the same amount). The second difference equals $2 \times$ the coefficient of $n^{2}$.

\textbf{Worked example.} Find the $n$th term of $2, 5, 10, 17, \dots$

First differences are $3, 5, 7$; the second difference is $2$, so the $n^{2}$ part is $1n^{2}$. Subtracting $n^{2}$ ($1, 4, 9, 16$) from the sequence leaves $1, 1, 1, 1$. So the $n$th term is $n^{2} + 1$.

A \textbf{cubic} 三次 sequence such as $1, 8, 27, 64, \dots$ has $n$th term $n^{3}$. An \textbf{exponential sequence} 指数数列 such as $2, 6, 18, 54, \dots$ multiplies by a fixed number each time; here the $n$th term is $2 \times 3^{\,n-1}$.

\section*{Direct and inverse proportion (Extended)}

Two quantities are in \textbf{direct proportion} 正比例 if one is always a fixed multiple of the other: $y \propto x$ means $y = kx$, where $k$ is a \textbf{constant} 常数. They are in \textbf{inverse proportion} 反比例 if one rises as the other falls: $y \propto \dfrac{1}{x}$ means $y = \dfrac{k}{x}$. The symbol $\propto$ is read "is proportional to". You can also have proportion to a square, square root, cube or cube root.

\textbf{Worked example.} $y$ is in direct proportion to $x$, and $y = 12$ when $x = 3$. Find $y$ when $x = 7$.

First find $k$: $12 = k \times 3$, so $k = 4$ and $y = 4x$. Then $y = 4 \times 7 = 28$.

\section*{Graphs in practical situations}

The \textbf{gradient} 斜率 (steepness) of a graph shows a \textbf{rate of change} 变化率.

\begin{itemize}
\item A \textbf{travel graph} 行程图 (distance–time graph) has gradient equal to speed; a flat part means the object is not moving.

\item A \textbf{conversion graph} 换算图 is a straight line used to change between two units (for example, miles and kilometres).

\end{itemize}

\subsection*{Speed–time graphs (Extended)}

On a speed–time graph the gradient is the \textbf{acceleration} 加速度 (or \textbf{deceleration} 减速度 if the speed falls), and the \textbf{area} 面积 under the graph is the \textbf{distance} 距离 travelled.

\textbf{Worked example.} A car speeds up from rest to $20\text{ m/s}$ in $8\text{ s}$, then stays at $20\text{ m/s}$ for $12\text{ s}$. Find the acceleration and the total distance.

Acceleration $= \dfrac{20}{8} = 2.5\text{ m/s}^{2}$. The distance is the area: a triangle plus a rectangle,

\[
\tfrac{1}{2} \times 8 \times 20 + 12 \times 20 = 80 + 240 = 320\text{ m}.
\]

\section*{Graphs of functions}

To draw a graph, make a \textbf{table of values} 数值表: choose values of $x$, work out $y$, then plot the \textbf{coordinates} 坐标 and join them with a smooth curve.

The points where a graph crosses the $x$-axis (the horizontal \textbf{axis} 坐标轴) are the \textbf{roots} 根 — the solutions of $y = 0$.

You can solve an equation by reading a graph. The \textbf{intersection point} 交点 of a line and a curve gives the solution of the two equations together.

For \textbf{exponential growth} 指数增长 and \textbf{exponential decay} 指数衰减, the graph of $y = a\,b^{x} + c$ rises (or falls) faster and faster and flattens towards a horizontal line.

\section*{Sketching curves}

A quick sketch should show the right shape and the key features: where it crosses the axes, any \textbf{symmetry} 对称, and any line the curve gets close to.

\begin{center}
\adjustbox{max width=\linewidth}{%
\begin{tabular}{ll}
\toprule
\textbf{Function} & \textbf{Shape} \\
\midrule
linear, $y = mx + c$ & straight line; gradient $m$, $y$-\textbf{intercept} 截距 $c$ \\
quadratic, $y = ax^{2} + bx + c$ & a \textbf{parabola} 抛物线 (U-shape if $a>0$, $\cap$-shape if $a<0$) \\
cubic, $y = ax^{3} + bx + c$ & an S-shaped curve \\
reciprocal, $y = \dfrac{a}{x} + b$ & two separate curves \\
exponential, $y = a\,r^{x} + b$ & fast growth or decay \\
\bottomrule
\end{tabular}}
\end{center}

For a parabola, completing the square gives the \textbf{turning point} 转折点 (the lowest or highest point). For example $y = (x + 3)^{2} - 8$ has its turning point at $(-3, -8)$.

An \textbf{asymptote} 渐近线 is a line that the curve gets closer and closer to but never touches — for example, the $x$-axis for $y = \dfrac{a}{x}$, or the line $y = b$ for $y = a\,r^{x} + b$.

\section*{Differentiation (Extended)}

\textbf{Differentiation} 微分 finds the gradient of a curve at any point. You can estimate it by drawing a \textbf{tangent} 切线 (a line that just touches the curve) and measuring its gradient.

The exact rule: if $y = ax^{n}$, then the \textbf{derivative} 导数 is

\[
\frac{\mathrm{d}y}{\mathrm{d}x} = a\,n\,x^{\,n-1}.
\]

Differentiate a sum term by term.

\textbf{Worked example.} If $y = x^{3} + 2x^{2} - 5x$, then $\dfrac{\mathrm{d}y}{\mathrm{d}x} = 3x^{2} + 4x - 5$.

A \textbf{stationary point} 驻点 (turning point) is where the gradient is zero, so set $\dfrac{\mathrm{d}y}{\mathrm{d}x} = 0$.

\textbf{Worked example.} Find the turning point of $y = x^{2} - 6x + 5$.

\[
\frac{\mathrm{d}y}{\mathrm{d}x} = 2x - 6 = 0 \;\Rightarrow\; x = 3, \quad y = 3^{2} - 6(3) + 5 = -4.
\]

The turning point is $(3, -4)$. To decide whether a turning point is a \textbf{maximum} 最大值 or a \textbf{minimum} 最小值, check the sign of the gradient on each side, or use the second derivative (positive means a minimum).

\section*{Functions (Extended)}

A \textbf{function} 函数 turns each input into one output. We write $f(x)$, for example $f(x) = 3x - 5$, so $f(2) = 1$. The set of allowed inputs is the \textbf{domain} 定义域; the set of possible outputs is the \textbf{range} 值域.

The \textbf{inverse function} 反函数 $f^{-1}(x)$ undoes the function. To find it, write $y = f(x)$, swap the roles, and make $x$ the subject.

\textbf{Worked example.} Find the inverse of $f(x) = 3x - 5$.

\[
y = 3x - 5 \;\Rightarrow\; x = \frac{y + 5}{3}, \quad \text{so} \quad f^{-1}(x) = \frac{x + 5}{3}.
\]

A \textbf{composite function} 复合函数 applies one function after another: $gf(x)$ means "do $f$ first, then $g$".

\textbf{Worked example.} If $f(x) = 2x$ and $g(x) = x + 3$, then

\[
gf(x) = g(2x) = 2x + 3, \qquad fg(x) = f(x + 3) = 2(x + 3) = 2x + 6.
\]


\end{document}
